% SynAthlete-AIoT: IEEE conference draft (D-009). Open in Overleaf (IEEEtran class).
% Every number comes from generated/*.tex, written by `python tasks.py paper-figures` from the
% experiments/**/metrics.json files (AGENTS.md invariant 1). Do not type numbers by hand.
\documentclass[conference]{IEEEtran}
\usepackage{graphicx,booktabs,amsmath,url,xspace}
% generated by src/paper/figures.py: do not edit
\newcommand{\canaryAuc}{0.498 (day) and 0.513 (week)\xspace}
\newcommand{\explainSummary}{Explanation checks: n/a.\xspace}
\newcommand{\fusedAuc}{n/a\xspace}
\newcommand{\fusedPrAuc}{n/a\xspace}
\newcommand{\injuryAucRange}{0.58--0.63\xspace}
\newcommand{\labelFreeAuc}{0.584 (day) and 0.589 (week)\xspace}
\newcommand{\lockboxSummary}{Pending the evaluation freeze.\xspace}
\newcommand{\lovdalRunners}{74\xspace}
\newcommand{\pamapRealFone}{0.632\xspace}
\newcommand{\reproDayAuc}{0.617 (95\% CI 0.573--0.664)\xspace}
\newcommand{\tiersReduction}{n/a\xspace}


\begin{document}
\title{SynAthlete-AIoT: An Honest, Explainable AIoT Pipeline for Athlete Injury-Risk Prediction
with Synthetic Wearable Data}
\author{\IEEEauthorblockN{[Author names: to be added]}
\IEEEauthorblockA{[Department], SRM Institute of Science and Technology, India}}
\maketitle

\begin{abstract}
Wearable artificial-intelligence-of-things (AIoT) systems promise early warning of athlete
injuries, but published frameworks often stop at the design, and reported accuracies can be
inflated by evaluation shortcuts. We implement a complete pipeline (synthetic wearable data,
simulated edge--fog--cloud tiers, a fused sequence-plus-tabular injury model, checked
explanations and a coach dashboard) and evaluate it under a pre-registered, person-level
protocol on real injury data from 74 competitive runners. Learned models predict injuries only
modestly above chance (ROC-AUC \injuryAucRange{} over five development groups); a faithful
rebuild of a published bagged-XGBoost method reproduces its behaviour, and without its
label-reading normalisation it performs like simple baselines. Edge processing sends
\tiersReduction{} less data to the cloud. We report every model, failures included, and release
the protocol and code.
\end{abstract}

\section{Introduction}
Coaches need to see injury risk building up over days, not in single snapshots. Recent AIoT
designs for athlete health propose wearables with edge--fog--cloud processing but leave the
prediction and explanation layers unimplemented~\cite{aiotframework2026}. Public data make the
problem harder: no public dataset pairs high-frequency physiological streams with real injury
outcomes. Daily training logs with real injuries do exist~\cite{lovdal2021,lovdal2021data},
while wearable datasets with rich physiology have no injury labels~\cite{reiss2012,schmidt2018}.

A second problem is evaluation. Injury events are rare, models are tuned on small cohorts, and
preprocessing can quietly use test labels; such leakage inflates results~\cite{kapoor2023}, and
model selection outside the test loop biases them~\cite{varma2006,cawley2010}.

\textbf{Contributions.} (1) An end-to-end, open AIoT pipeline in which each component has an
explicit, honest role. (2) A pre-registered, person-level evaluation protocol with a sealed
final test, a leak alarm and runner-level uncertainty. (3) A faithful reproduction of a
published injury-prediction method, with and without its label-reading normalisation.
(4) A generator comparison (TimeGAN, WGAN-GP, TimeVAE, classic augmentation) judged against
criteria fixed before any run. (5) Explanations checked by deletion, randomisation and stability
tests rather than assumed.

\section{Related Work}
\textit{Injury prediction.} L\"ovdal et al.\ predicted injuries of competitive runners from
daily logs with bagged XGBoost~\cite{lovdal2021,chen2016}. Tree ensembles remain strong on
tabular data~\cite{grinsztajn2022,breiman2001,prokhorenkova2018}; recurrent networks read days
in order~\cite{cho2014}. \textit{Synthetic time series.} TimeGAN~\cite{yoon2019},
WGAN-GP~\cite{gulrajani2017} and TimeVAE~\cite{desai2021} generate sequences; classic
augmentation~\cite{um2017} is a strong baseline. \textit{Explanations.} SHAP~\cite{lundberg2017}
and integrated gradients~\cite{sundararajan2017} attribute predictions to inputs; sanity checks
show that attributions must be tested~\cite{adebayo2018}. \textit{Evaluation.} With rare events,
PR-AUC is more informative than ROC-AUC~\cite{saito2015,davis2006}; calibration and decision
curves complete the picture~\cite{vancalster2019,vickers2006}.

\section{Data and Roles}
Table~\ref{tab:data} lists every dataset and its single role. Injury claims come only from the
L\"ovdal runner logs~\cite{lovdal2021data}: \lovdalRunners{} runners, a day file (7 days
$\times$ 10 features) and a week file (3 weeks $\times$ 22 features plus 3 relative changes).
PAMAP2~\cite{reiss2012} provides the wearable stream (heart rate and three accelerometers, 25~Hz,
5.12-s windows); WESAD~\cite{schmidt2018} tests generator generalisation; Sport DB 2.0 checks
heart-rate realism; SoccerMon~\cite{midoglu2024} and PMData are exploratory. Datasets of unknown
origin appear only as labelled demos. Synthetic data is never evidence about real injuries.
\begin{table}[t]\centering\caption{Datasets and their single role}\label{tab:data}
\footnotesize\begin{tabular}{p{2.6cm}p{1.2cm}p{3.6cm}}\toprule
Dataset & Origin & Role \\\midrule
Lövdal runners (74 runners, 7 years of daily logs) & real & the only injury evidence: baselines, combined model, explanations, coach view \\
Kaggle 'Injury Prediction for Competitive Runners' & real (copy) & byte-identical copy of Lövdal: no separate role \\
PAMAP2 (9 people, wearable sensors) & real & edge sensor stream, activity recognition, generator training and fidelity \\
WESAD (15 people, wrist sensors) & real & second dataset for generator generalisation \\
Sport DB 2.0 (55 heart-rate recordings) & real & heart-rate realism check and second edge stream \\
SoccerMon, PMData & real & exploratory only (Track C), never main evidence \\
SIRP-600 & unknown origin & demo only: shows the SHAP code runs \\
ziya07 wearable & unknown origin & demo only: dashboard demo stream \\
\bottomrule\end{tabular}\end{table}


\section{Evaluation Protocol}
\textit{Splits.} All splits are by person. The ten most recently joined runners form a sealed
final test, opened once after an evaluation freeze. The other 64 runners form five development
rounds; in each, settings, Platt calibration, the alert level and early stopping use that
round's validation runners only, and its test runners are scored once.
\textit{Scoring.} The headline is PR-AUC (average precision) next to the injury rate (its
chance level), with ROC-AUC for comparison with the literature. The alert level minimises
$|\mathrm{sensitivity}-\mathrm{specificity}|$ on validation runners. Calibration is reported
after Platt scaling~\cite{vancalster2019}.
\textit{Tuning.} Eight settings per model and round by seeded random search~\cite{bergstra2012},
written before the first run; the chosen setting is refit with three seeds.
\textit{Leakage controls.} One z-score scaler fitted on training runners only; no identity,
date or own-history inputs (in this dataset, a runner's own event history reveals the label);
a shuffled-label alarm~\cite{ojala2010}; an independent code review.
\textit{Uncertainty.} Runner-level cluster bootstrap within test groups~\cite{field2007}.

\section{Methods}
\textit{Baselines.} Constant risk, an uncoupled ACWR-style rule~\cite{lolli2019}, logistic
regression, random forest, CatBoost and a GRU; the L\"ovdal bagged XGBoost as an exact copy
(reproduction check on stand-in runners) and as a label-free variant.
\textit{Fused model.} A sequence arm (GRU, LSTM or TCN~\cite{bai2018}) on day sequences and a
tabular arm (random forest or CatBoost) on week features, chosen per round on validation runners,
combined by logistic regression trained on out-of-fold arm predictions grouped by runner. The
calibrated output is shown as a 0--100 risk score (its percentile among validation events).
\textit{Generators.} Class-conditional TimeGAN, WGAN-GP and TimeVAE, trained only on each round's
training people, with settings chosen by validation train-on-synthetic-test-on-real (TSTR)
macro-F1; plus classic augmentation.
\textit{Tiers.} Edge (reflex check and activity CNN per window), fog (session summaries) and
cloud (daily injury model) run as separate processes; we compare bytes sent and decision latency
with a cloud-only design.
\textit{Explanations.} TreeSHAP on the tabular arm, integrated gradients on the sequence arm, an
explanation-only attention model shown next to integrated gradients, and 1D Grad-CAM on the
activity CNN; all checked by deletion, model-randomisation and stability tests with thresholds
fixed in advance.

\section{Results}
\begin{table}[t]\centering\caption{Injury prediction, mean $\pm$ SD over five development rounds (chance PR-AUC $\approx$ injury rate)}\label{tab:injury}
\footnotesize\begin{tabular}{llcc}\toprule
Model & File & PR-AUC & ROC-AUC \\\midrule
Constant risk & day & 0.013 $\pm$ 0.000 & 0.500 $\pm$ 0.000 \\
Constant risk & week & 0.013 $\pm$ 0.000 & 0.500 $\pm$ 0.000 \\
ACWR-style rule & week & 0.015 $\pm$ 0.001 & 0.540 $\pm$ 0.021 \\
Logistic regression & day & 0.022 $\pm$ 0.004 & 0.626 $\pm$ 0.019 \\
Logistic regression & week & 0.020 $\pm$ 0.004 & 0.621 $\pm$ 0.029 \\
Random forest & day & 0.019 $\pm$ 0.002 & 0.603 $\pm$ 0.024 \\
Random forest & week & 0.022 $\pm$ 0.007 & 0.614 $\pm$ 0.042 \\
CatBoost & day & 0.021 $\pm$ 0.004 & 0.598 $\pm$ 0.043 \\
CatBoost & week & 0.019 $\pm$ 0.002 & 0.606 $\pm$ 0.028 \\
GRU & day & 0.023 $\pm$ 0.006 & 0.631 $\pm$ 0.009 \\
L\"ovdal, label-free & day & 0.017 $\pm$ 0.003 & 0.584 $\pm$ 0.044 \\
L\"ovdal, label-free & week & 0.021 $\pm$ 0.006 & 0.589 $\pm$ 0.066 \\
\textbf{Fused} & shared & 0.022 $\pm$ 0.003 & 0.609 $\pm$ 0.024 \\
\bottomrule\end{tabular}\end{table}

\textit{Injury prediction.} Table~\ref{tab:injury} shows every model. The leak alarm passed
(shuffled-label ROC-AUC \canaryAuc{}). Learned models reach ROC-AUC \injuryAucRange{}, with
differences smaller than the spread between rounds. The fused model reaches ROC-AUC
\fusedAuc{} and PR-AUC \fusedPrAuc{}. The exact copy of the published method passed its
pre-registered check (day-file ROC-AUC \reproDayAuc{}); without its label-reading normalisation
the same model reaches ROC-AUC \labelFreeAuc{}.

\begin{table}[t]\centering\caption{Generators on PAMAP2 (4 rounds): discriminative accuracy (0.5 = indistinguishable), TSTR and real+synthetic macro-F1, pre-registered criteria passed}\label{tab:gen}
\footnotesize\begin{tabular}{lcccc}\toprule
Generator & Disc.\ acc. & TSTR & Real+syn. & Passed \\\midrule
\multicolumn{5}{c}{n/a} \\
\bottomrule\end{tabular}\end{table}

\textit{Synthetic data.} Table~\ref{tab:gen} compares the generators with real-data training
(real-only macro-F1 \pamapRealFone{}).

\begin{table}[t]\centering\caption{Edge--fog--cloud vs cloud-only (simulated on one computer)}\label{tab:tiers}
\footnotesize\begin{tabular}{llrr}\toprule
Stream & Mode & Bytes into cloud & Median decision (ms) \\\midrule
\multicolumn{4}{c}{n/a} \\
\bottomrule\end{tabular}\end{table}

\textit{Tiers.} Processing at the edge reduced the data sent to the cloud by \tiersReduction{}
(Table~\ref{tab:tiers}).

\textit{Explanations.} \explainSummary{}

\textit{Final test.} \lockboxSummary{}

\section{Discussion and Limitations}
Training logs alone carry little warning before an injury; honest models are only modestly
better than chance, and selection on test data or label-reading normalisation would have hidden
this. The dataset's authors removed healthy days close to injuries, so alert rates are not daily
rates and net benefit is not clinical utility. The tiers are simulated on one computer, and the
wearable streams come from different people than the runners: they demonstrate the system, not
injury evidence. Synthetic data was used for training only.

\section{Conclusion}
We implemented and honestly evaluated a complete explainable AIoT pipeline for injury risk.
The protocol, code and all results, including negative ones, are released for reuse.

\section*{AI-Use Statement}
AI coding assistants (Claude Code; Antigravity with Claude Opus 4.6 and Gemini 3.8 Flash) helped
write code and text under human direction; every AI contribution is logged in the project's
AI\_LOG, every number is produced by code, and all claims were checked by the authors.

\bibliographystyle{IEEEtran}
\bibliography{refs}
\end{document}
